\begin{abstract}
\noindent
We re-examine the relation between quarterly real GDP growth and changes in the
unemployment rate in the \val{n_countries} member states of the European Union from 2009
to 2026, using official Eurostat data frozen with cryptographic hashes and analyses
registered before estimation.

\textbf{Association.} With 2020--2021 excluded, the two-lag Okun coefficient sum is
negative, with a 95\% interval below zero, in \val{cl_n_neg_excluded} of
\val{n_countries} countries. Including the pandemic quarters pulls most sums towards zero.
A hierarchical Bayesian model with Student-$t$ errors, residual autocorrelation and a
common European factor gives an unweighted cross-country mean of \val{by_mu} (90\%
interval \val{by_mu_ci}), with a between-country standard deviation of \val{by_tau}.
Country ranks remain highly uncertain.

\textbf{Prediction.} Hierarchical pooling does not improve one-quarter-ahead forecasts
beyond sampling variation, and neither do lagged GDP growth or the common factor: none of
the pre-registered forecast comparisons is decided at the 5\% familywise level.

\textbf{Causation.} External instruments for monetary-policy and oil-supply shocks are weak
for real activity, with first-stage $|t|$ at most \val{id_max_t} quarterly and
\val{mo_max_t} monthly. A foreign-demand shift-share instrument is infeasible (expected
first-stage $t$ of \val{ss_P1_t}). Sign restrictions that the data admit leave the
shock-conditional Okun ratios unbounded.

\textbf{Countries.} We compare the results with three earlier claims: a large association
for Spain, a delayed one for the Netherlands and a weak one for Luxembourg. The Spanish claim
is supported in every accepted specification. The Dutch and Luxembourg claims are not
resolved by their registered rules. In Luxembourg, the data give no sign that the
cross-border measurement mismatch explains the weak association: domestic and resident
employment barely move with GDP, and output per hour absorbs the fluctuations.

Every Okun coefficient in this study is a descriptive association; none is an identified
causal effect.
\end{abstract}
